\subsection{The integers modulo $n$}\label{sec:modn}

This is the first group in the gallery that is not presented as the symmetries of something you can
hold, and the first whose multiplication is given by a \emph{formula with cases}. Both facts mean
the axioms have to be checked by hand, and the third one costs real work.

\emph{The set.} $\Z/n\Z=\{0,1,\dots,n-1\}$, for a fixed integer $n>1$.

\emph{The rule.} Add, and if the answer reaches $n$, subtract $n$:
\[a\oplus b=\begin{cases} a+b, & a+b\le n-1,\\ a+b-n, & a+b\ge n.\end{cases}\]

\subsubsection*{What the elements really are}

Each symbol $0,\dots,n-1$ is a name for a whole family of integers: those leaving that remainder on
division by $n$. Colour every integer by its remainder and the picture is a repeating pattern; wrap
the line round a circle with $n$ points and each colour lands on one point.

\sfig{Why there are exactly five elements when $n=5$. Same colour means same remainder, hence the
same element of $\Z/5\Z$; the integers $\dots,-5,0,5,10,\dots$ all go by the name $0$.}{\figXad}

\begin{why}
\emph{Why the rule needs only two cases.} Both $a$ and $b$ are at most $n-1$, so the ordinary sum
$a+b$ is at most $2n-2$, which is less than $2n$. One subtraction therefore always suffices, and
$a\oplus b$ is exactly the remainder of $a+b$ on division by $n$. This is the reason the operation
is called \emph{arithmetic modulo $n$} and written $a+b\equiv c \pmod n$.
\end{why}

\sfigs{0.90}{The rule on the circle for $n=5$. Left: no wrap. Middle: the walk passes $0$, so $n$ is
subtracted. Right: the inverse of $2$ is $3$, because three more steps complete a lap.}{\figXae}

\subsubsection*{The three axioms}

\emph{Identity.} $0$ works: $a+0=a\le n-1$, so no wrap occurs and $0\oplus a=a\oplus 0=a$. On the
circle, adding $0$ is standing still.

\emph{Inverses.} The inverse of $0$ is $0$. For $m\neq 0$ the inverse is $n-m$, because
$m+(n-m)=n\ge n$, so the rule subtracts $n$ and the answer is $0$. On the circle, walking $n-m$
steps on from $m$ completes exactly one lap.

\emph{Associativity.} Here the cases bite: each $\oplus$ may or may not subtract $n$, so comparing
$(a\oplus b)\oplus c$ with $a\oplus(b\oplus c)$ seems to need a list of combinations. A worked case
with $n=5$, $a=3$, $b=4$, $c=2$:

\begin{center}
\small
\begin{tabular}{@{}ll@{}}
\toprule
$(a\oplus b)\oplus c$ & $a\oplus(b\oplus c)$\\
\midrule
$3\oplus4=7-5=2$ & $4\oplus2=6-5=1$\\
$2\oplus2=4$, no wrap & $3\oplus1=4$, no wrap\\
\midrule
\multicolumn{2}{c}{both equal $4$, the remainder of $3+4+2=9$ on division by $5$}\\
\bottomrule
\end{tabular}
\end{center}

\begin{why}
\emph{The argument, with no cases at all.} Each application of $\oplus$ either keeps the sum or
subtracts $n$. So however you bracket, the answer is the integer $a+b+c$ minus some multiple of
$n$; and it lies in $\{0,\dots,n-1\}$, because that is where $\oplus$ always lands. But there is
only \emph{one} integer in $\{0,\dots,n-1\}$ differing from $a+b+c$ by a multiple of $n$, namely its
remainder. Both bracketings therefore equal that remainder, hence each other. In one sentence:
$\oplus$ is ``add in $\Z$, then take the remainder'', and taking remainders is compatible with
addition.
\end{why}

\begin{pitfall}
The word ``multiplication'' in Definition~\ref{def:group} here means \emph{addition}. Ordinary
multiplication modulo $n$ is not a group rule on this set, because $0$ has no multiplicative
inverse. Note also that $\Z/n\Z$ is not a subset of $\Z$ in any useful sense: $3\oplus 4=2$ in
$\Z/5\Z$ while $3+4=7$ in $\Z$.
\end{pitfall}

\subsubsection*{The extra properties}

$|\Z/n\Z|=n$, and the group is \defn{abelian}, since $a+b=b+a$ in $\Z$ and the wrap does not care
about order. It is \defn{cyclic}: starting at $0$ and adding $1$ repeatedly visits every element, so
$\{1\}$ generates it all by itself. More generally $m$ generates $\Z/n\Z$ exactly when $m$ and $n$
have no common factor.

\begin{why}
\emph{Why the condition is $\gcd(m,n)=1$.} If $m$ and $n$ have no common factor, there are integers
$u,v$ with $um+vn=1$; reducing modulo $n$ gives $u\cdot m\equiv 1$, so adding $m$ to itself $u$
times produces $1$, and $1$ generates everything. Conversely, if $d>1$ divides both $m$ and $n$,
then every multiple of $m$ is again a multiple of $d$ modulo $n$, so $1$ is never reached.
\end{why}

\sfig{Familiar instances: the $12$-hour clock is $\Z/12\Z$, and $9+4=13$ is renamed $1$.}{\figXaf}

\sfig{Subgroups of $\Z/12\Z$. Each highlighted set is closed under ``add and wrap'', so each is a
group in its own right, and each is a regular polygon inscribed in the clock, one for every divisor
of $12$. Since $5$ and $12$ have no common factor, adding $5$ repeatedly visits all twelve points;
starting from $4$ you only ever reach the triangle.}{%
\figclocksub{6}{$\{0,6\}$\\generated by $6$}\hspace{4pt}%
\figclocksub{4}{$\{0,4,8\}$\\generated by $4$}\hspace{4pt}%
\figclocksub{3}{$\{0,3,6,9\}$\\generated by $3$}\hspace{4pt}%
\figclocksub{2}{$\{0,2,\dots,10\}$\\generated by $2$}}

\begin{tryit}
In $\Z/12\Z$: what is the inverse of $5$? of $8$? which elements are their own inverses? Check
associativity for $a=7$, $b=9$, $c=8$ by both bracketings, and against the remainder of $24$. Then
say which axiom fails for the rule ``subtract, then wrap''.
\end{tryit}

\subsection{Decorating a polygon to see $\Z/n\Z$}\label{sec:arrowgon}

$\Z/n\Z$ was defined by a formula, but it is a symmetry group after all, and the object it is the
symmetry group of is obtained from the $n$-gon by a trick worth learning, because the same trick
will be used three more times.

The symmetries of a regular $n$-gon form $D_n$, which is twice too big. So spoil the polygon: draw
a little arrow on each edge, all running the same way round. A rotation carries arrows to arrows
pointing the same way round and is still a symmetry of the decorated polygon. A reflection reverses
every arrow, and is not. The decoration has removed exactly the reflections.

\sfigs{0.90}{A rotation preserves the arrows; a reflection reverses them. So the symmetry group of the
arrowed $n$-gon is the $n$ rotations, and the rotations compose by adding clicks modulo
$n$.}{\figXag}

\sfigs{0.90}{Composing clicks is addition modulo $n$: two clicks then four clicks is six clicks, and six
clicks of a pentagon is one click, because five clicks is a full turn. In symbols
$2+4\equiv1 \pmod 5$.}{\figXah}

\sfig{\label{fig:arrowgon}The arrowed $n$-gon, whose symmetry group is $\Z/n\Z$. It is also the Cayley graph of $\Z/n\Z$ with the generator $1$, in the sense of Section~\ref{sec:cayley}.}{\figarrowgon}

\begin{bigpic}
This is the first example of the chapter's main theme in miniature. An algebraic object, $\Z/n\Z$,
was matched with a geometric one, the arrowed polygon, and the matching turned addition with wrap
into rotation. Section~\ref{sec:cayley} shows this is not a lucky accident: \emph{every} group can
be presented this way, once you know how to build the right decorated object.
\end{bigpic}

\subsection{The integers as slides of a line}\label{sec:Z}

\emph{The set.} $\Z$. \emph{The rule.} Addition. Identity $0$, inverse $-n$, associativity
inherited from ordinary arithmetic. So $\Z$ is a group, and an abelian one.

To see it as a symmetry group, mark the integer points on a line. A rigid motion of the line that
sends marked points to marked points is either a slide by a whole number of units, or a reflection
about a marked point, or a reflection about a midpoint between two neighbours. That collection is
bigger than $\Z$; it is called the \defn{infinite dihedral group} $D_\infty$, the $n\to\infty$
version of $D_n$.



\begin{why}
\emph{Why the slide passes both tests and the reflection does not.} Sliding by $3$ is the map
$x\mapsto x+3$, the same shift at every point, so gaps are preserved,
$|(x+3)-(y+3)|=|x-y|$, and order is preserved, $x<y\Rightarrow x+3<y+3$. An arrow records only
``$k$ is to the left of $k+1$'', so it keeps pointing right. Reflecting about the point $2$ is
$x\mapsto 4-x$: gaps are preserved too, but now $x<y$ becomes $4-x>4-y$, so every arrow turns round.
\end{why}

\sfigs{0.90}{\label{fig:arrowline}Composing slides is adding integers: slide by $+3$, then by $-2$, and the net effect is $+1$. The arrowed line is the Cayley graph of $\Z$.}{\figXak}

\sfig{Subgroups of $\Z$. Every subgroup is the set of multiples of some $m\ge0$. The map
$n\mapsto 2n$ shows $\Z$ containing a smaller copy of itself, something no finite group can
do.}{\figsubZ}

\subsection{The plane lattice $\Z^2$}\label{sec:Z2}

\emph{The set.} $\Z^2=\{(m,n): m,n\in\Z\}$, pairs of integers. \emph{The rule.} Add coordinates:
$(m,n)+(m',n')=(m+m',n+n')$. Identity $(0,0)$, inverse $(-m,-n)$, associativity inherited
coordinatewise. Abelian.

Geometrically, $(m,n)$ is the slide ``$m$ to the right and $n$ up'', and composing two slides adds
the coordinates, which is vector addition.

\sfigs{0.98}{\label{fig:slides}Slides of the plane compose by adding coordinates. Every point of the
lattice moves by the same amount, so it is enough to follow one of them: the ringed red point. The
middle panel also breaks the slide $(2,1)$ into unit steps right, right and up, so
$(2,1)=2(1,0)+1(0,1)$ and the two unit slides generate every slide.}{\figSlide}

The object whose symmetries are exactly $\Z^2$ is the integer grid, decorated so that rotations and
reflections are spoiled: colour the horizontal edges one colour and the vertical edges another.

\sfigs{0.99}{\label{fig:gridtest}Which motions of the plane survive the decoration. The background
grid is the same in all three panels: it is the object. Follow the highlighted edge. A slide carries
it to another horizontal edge, so blue still meets blue and the slide is a symmetry. A quarter turn
carries it to a vertical edge, where a blue arrow would have to sit on a green one, so the quarter
turn is not. A reflection fails for a different reason: it reverses the arrows. Only the slides
survive, and they compose exactly like $\Z^2$.}{\figGridCheck}

\sfigs{0.62}{\label{fig:grid}The decorated grid, whose symmetry group is $\Z^2$, and which is also its Cayley graph with the two unit slides as generators.}{\figlattice}

\begin{pitfall}
$\Z^2$ is abelian, and it is generated by two elements, $(1,0)$ and $(0,1)$. In
Section~\ref{sec:free} we meet another group with two generators which is as far from abelian as
possible. Two generators say very little about a group on their own.
\end{pitfall}

\subsection{Matrix groups, in one paragraph}\label{sec:matrix}

For completeness, one important family that is not a symmetry group of a polygon or a cube.
Let $\GL(n,\R)$ be the set of $n\times n$ matrices with real entries and nonzero determinant, with
matrix multiplication as the rule. Closure holds because $\det(AB)=\det(A)\det(B)$ is nonzero when
both factors are; the identity is the identity matrix; inverses exist exactly because the
determinant is nonzero; and associativity holds because matrix multiplication represents
composition of linear maps, which is composition of functions. Requiring determinant exactly $1$
gives a subgroup $\SL(n,\R)$. Replacing $\R$ by $\Z/p\Z$ for a prime $p$ gives finite examples: the
invertible $2\times2$ matrices over $\Z/2\Z$ form a group with $6$ elements which behaves exactly
like $S_3$, since it permutes the three nonzero vectors of the plane over $\Z/2\Z$.
